\chapter*{Part Synthesis: Learning from Interaction}
\addcontentsline{toc}{chapter}{Part Synthesis: Learning from Interaction}

This short part isolates a category change. Until now, the book has trained models that fit a function to a static dataset; here, the system must \emph{act}, observe consequences, and revise its behavior. Reinforcement learning deserves its own part because reward design replaces label design, exploration joins evaluation as a first-class design constraint, and the data distribution is shaped by the agent's own policy rather than fixed in advance.

\begin{chaptersummary}
\item \textbf{Question this part taught.} When does the learning problem stop being ``fit a function to data'' and start being ``decide what to do next, observe what happened, and update''?
\item \textbf{Artifact chain this part added.} Interaction Design Card---a structured record of state, action space, reward, exploration strategy, evaluation plan, and guardrails before any RL system is treated as deployable.
\item \textbf{What a student can now do.} Recognize when a problem genuinely requires interaction rather than prediction, formalize it as an MDP, choose between value-based, policy-based, and actor-critic methods, audit a reward signal for misalignment, and document the safety and exploration assumptions that decide whether the system can be tested in the wild.
\item \textbf{What comes next.} Part~IV returns to predictive learning, but with the interaction lens still active: vision, language, sequences, ranked recommendation, and tabular pipelines all involve outputs, row definitions, or operating policies that change what the system later sees, and the reward-design and exposure-bias lessons from this part keep mattering even when the surface-level task looks supervised.
\end{chaptersummary}

The course-support assistant and the campus shuttle reappear here in their action-taking forms: the assistant must decide which response actually resolves a student's issue, and the shuttle must decide when to brake. Both make the categorical shift visible---a fluent prediction is not the same as a good action, and a labeled dataset is not the same as a reward signal. Part~IV picks the supervised thread back up across modalities, including tabular row-level decisions, while carrying forward the interaction discipline introduced here.
